% ============================================================
% body.tex —— 答案册·P1 轮4 S5【组装器产出】——勿手改（组装S5.py 为唯一值源，防回卷）
% 生成：2026-09-12｜现值键：208｜缺账占位：0｜PENDING 键：0｜对号契约：每键一行 `%% pair:键`
% 版式：册首行/章标/节标/分区组行/条目宏在 main.tex＋qp-answ-*（PL1 冻结原样）；悬挂两档＝单号 5.8mm／双位 7.6mm。
% 收尾轮（2026-09-12）：批5b 65 键回填转实值（白名单撤项、钉-8/钉-9 哨兵转生效、钉-10～15 六值门生效）；导-课时92/93-预习填空随批2/批3 §二-A1 补登回填。
% ============================================================
\setcounter{page}{1}   % 独立册起页 1；装配轮按四本跨本连续页码改写
\par\nointerlineskip
{\centering\fontsize{16.88pt}{22pt}\selectfont\hejie 参考答案\par}
\vspace{2.4mm}
\zhangtitle{第九章\quad 静电场及其应用}

% ============================================================
% 【甲】导学件答案（4 节：课时91~94＋章末·本章易错过关 5 题）——整册单 multicols
% （节标入栏内，照数学 M2 【甲】同构；栏断由 multicol 自理，页 OTR 不背块序；
%   唯一手动栏断＝章末前 \columnbreak，收尾轮消 9.4 尾行溢栏造成的单行空页）
% ============================================================
\begin{multicols}{2}
\raggedcolumns
\emergencystretch=1em
\jietitle{9.1\quad 电荷(导学件答案)}
\qufen{第1课时}{课前预习（填空＋判断6）＋课中探究（例1×3／变式1×3／拓展延伸）＋课堂评价5}
% ---- 课时91（同构槽序：预习填空→诊断判断6→探究点一二三→拓展延伸→课堂评价5）----
% pair:导-课时91-预习填空
\ansline{预习填空}{知识点一：正；负；排斥；吸引；电子；负；正；平分；中和；平分。知识点二：张开；带；闭合；中和；自由电子；负；正。知识点三：创生；保持不变；最小；$1.60\times10^{-19}$；整数倍。}
\vskip 0pt minus 1.2pt
% pair:导-课时91-判1
% pair:导-课时91-判2
% pair:导-课时91-判3
% pair:导-课时91-判4
% pair:导-课时91-判5
% pair:导-课时91-判6
\ansline{答案}{(1)$\times$\quad (2)\(\surd\)\quad (3)\(\surd\)\quad (4)$\times$\quad (5)\(\surd\)\quad (6)$\times$}
\vskip 0pt minus 1.2pt
% pair:导-课时91-探1-例1
% pair:导-课时91-探1-变式1
\ansline{探究一}{例1 C；变式1 D}
\vskip 0pt minus 1.2pt
% pair:导-课时91-探2-例1
% pair:导-课时91-探2-变式1
\ansline{探究二}{例1 B；变式1 A 正／B 负；分开后 A 正；先移棒后分＝不带}
\vskip 0pt minus 1.2pt
\ansline{解析}{中3（定稿指定「导学位」取材件）详解 CD 段接地时序按主脑回执②展开：移 C 前先不断地线，电荷流回大地，导体不带电；若先断地线再移 C，则导体得异号净电荷（就绪报告 §七-2）。}
\vskip 0pt minus 1.2pt
% pair:导-课时91-探3-例1
% pair:导-课时91-探3-变式1
\ansline{探究三}{例1 (1) 各 $+2.0\times10^{-9}\,\unit{C}$；(2) 甲、丙各 $+1.0\times10^{-9}\,\unit{C}$；变式1 $+3.0\times10^{-9}\,\unit{C}$}
\vskip 0pt minus 1.2pt
\ansline{解析}{（(2) 中甲承 (1) 之结果）甲带 $+2.0\times10^{-9}\,\unit{C}$ 与不带电丙接触再分开、电荷平分，故甲、丙各 $+1.0\times10^{-9}\,\unit{C}$（非取题干原值 $+6.0\times10^{-9}\,\unit{C}$ 独立分配）。}
\vskip 0pt minus 1.2pt
% pair:导-课时91-拓展
\ansline{拓展延伸}{各 $-q$（$\dfrac{+q-3q}{2}=-q$；异种电荷接触后电性由代数和正负决定）}
\vskip 0pt minus 1.2pt
% pair:导-课时91-评1
% pair:导-课时91-评2
% pair:导-课时91-评3
% pair:导-课时91-评4
% pair:导-课时91-评5
\ansline{评价}{1．B；2．B；3．C；4．甲、丙各 $+2.0\times10^{-9}\,\unit{C}$；5．$3.0\times10^{10}$ 个}
\vskip 0pt minus 1.2pt
\jietitle{9.2\quad 库仑定律(导学件答案)}
\qufen{第2课时}{课前预习（填空＋判断6）＋课中探究（例1×3／变式1×3／拓展延伸）＋课堂评价5}
% ---- 课时92（同构槽序：预习填空→诊断判断6→探究点一二三→拓展延伸→课堂评价5）----
% pair:导-课时92-预习填空
\ansline{预习填空}{知识点一：点电荷；理想化模型；正；反；连线；$k\dfrac{q_1q_2}{r^2}$；$9.0\times10^{9}$。知识点二：排斥；吸引；真空；静止；绝对值。知识点三：矢量和。}
\vskip 0pt minus 1.2pt
% pair:导-课时92-判1
% pair:导-课时92-判2
% pair:导-课时92-判3
% pair:导-课时92-判4
% pair:导-课时92-判5
% pair:导-课时92-判6
\ansline{答案}{(1)$\times$\quad (2)\(\surd\)\quad (3)\(\surd\)\quad (4)$\times$\quad (5)\(\surd\)\quad (6)$\times$}
\vskip 0pt minus 1.2pt
% pair:导-课时92-探1-例1
% pair:导-课时92-探1-变式1
\ansline{探究一}{例1 B；变式1 D（$4F$）}
\vskip 0pt minus 1.2pt
% pair:导-课时92-探2-例1
% pair:导-课时92-探2-变式1
\ansline{探究二}{例1 C（$\dfrac{4F}{3}$）；变式1 D（$0$）}
\vskip 0pt minus 1.2pt
% pair:导-课时92-探3-例1
% pair:导-课时92-探3-变式1
\ansline{探究三}{例1 $0.25\,\unit{N}$，方向沿另外两个点电荷连线的垂直平分线向外（教材例2 照用）；变式1 大小 $\dfrac{20kq^{2}}{L^{2}}$，方向沿 AB 连线指向 B}
\vskip 0pt minus 1.2pt
% pair:导-课时92-拓展
\ansline{拓展延伸}{$\dfrac{1}{4}$（$F\propto\dfrac{1}{r^{2}}$）}
\vskip 0pt minus 1.2pt
% pair:导-课时92-评1
% pair:导-课时92-评2
% pair:导-课时92-评3
% pair:导-课时92-评4
% pair:导-课时92-评5
\ansline{评价}{1．B；2．A（$7.2\times10^{-4}\,\unit{N}$）；3．B（靠近 A 一侧）；4．$\dfrac{F}{4}$；5．$\dfrac{F}{8}$}
\vskip 0pt minus 1.2pt
\jietitle{9.3\quad 电场 电场强度(导学件答案)}
\qufen{第3课时}{课前预习（填空＋判断6）＋课中探究（例1×3／变式1×3／拓展延伸）＋课堂评价5}
% ---- 课时93（同构槽序：预习填空→诊断判断6→探究点一二三→拓展延伸→课堂评价5）----
% pair:导-课时93-预习填空
\ansline{预习填空}{知识点一：电场；力；电荷量；$\dfrac{F}{q}$；牛顿每库仑；矢量；正；电场；无关。知识点二：$k\dfrac{Q}{r^{2}}$；背离；指向；矢量和。知识点三：切线；假想；大；相等；相同；平行直线。}
\vskip 0pt minus 1.2pt
% pair:导-课时93-判1
% pair:导-课时93-判2
% pair:导-课时93-判3
% pair:导-课时93-判4
% pair:导-课时93-判5
% pair:导-课时93-判6
\ansline{答案}{(1)\(\surd\)\quad (2)$\times$\quad (3)\(\surd\)\quad (4)\(\surd\)\quad (5)$\times$\quad (6)$\times$}
\vskip 0pt minus 1.2pt
% pair:导-课时93-探1-例1
% pair:导-课时93-探1-变式1
\ansline{探究一}{例1 C；变式1 $5.0\times10^{3}\,\unit{N/C}$，水平向右}
\vskip 0pt minus 1.2pt
% pair:导-课时93-探2-例1
% pair:导-课时93-探2-变式1
\ansline{探究二}{例1 $0$；变式1 $\dfrac{40kq}{L^{2}}$，由甲指向乙}
\vskip 0pt minus 1.2pt
% pair:导-课时93-探3-例1
% pair:导-课时93-探3-变式1
\ansline{探究三}{例1 B；变式1 C（保持不变）}
\vskip 0pt minus 1.2pt
% pair:导-课时93-拓展
\ansline{拓展延伸}{$0$（对径微元两两等大反向抵消）}
\vskip 0pt minus 1.2pt
% pair:导-课时93-评1
% pair:导-课时93-评2
% pair:导-课时93-评3
% pair:导-课时93-评4
% pair:导-课时93-评5
\ansline{评价}{1．A；2．C（$2E$）；3．B（$\dfrac{3kQ}{L^{2}}$，由 B 指向 A）；4．$1.0\times10^{-4}\,\unit{N}$，水平向右；5．$\dfrac{2kQ}{r^{2}}$，由正电荷指向负电荷}
\vskip 0pt minus 1.2pt
\jietitle{9.4\quad 静电的防止与利用(导学件答案)}
\qufen{第4课时}{课前预习（填空＋判断6）＋课中探究（例1×3／变式1×3／拓展延伸）＋课堂评价5}
% ---- 课时94（同构槽序：预习填空→诊断判断6→探究点一二三→拓展延伸→课堂评价5）----
% pair:导-课时94-预习填空
\ansline{预习填空}{知识点一：相反；处处为零；外表面；大；电离；尖端放电。知识点二：没有电荷；静电屏蔽；金属网；接地。知识点三：导走；接地；静电力。}
\vskip 0pt minus 1.2pt
% pair:导-课时94-判1
% pair:导-课时94-判2
% pair:导-课时94-判3
% pair:导-课时94-判4
% pair:导-课时94-判5
% pair:导-课时94-判6
\ansline{答案}{(1)\(\surd\)\quad (2)$\times$\quad (3)$\times$\quad (4)\(\surd\)\quad (5)\(\surd\)\quad (6)\(\surd\)}
\vskip 0pt minus 1.2pt
% pair:导-课时94-探1-例1
% pair:导-课时94-探1-变式1
\ansline{探究一}{例1 A；变式1 B}
\vskip 0pt minus 1.2pt
% pair:导-课时94-探2-例1
% pair:导-课时94-探2-变式1
\ansline{探究二}{例1 A；变式1 静电屏蔽；接地}
\vskip 0pt minus 1.2pt
% pair:导-课时94-探3-例1
% pair:导-课时94-探3-变式1
\ansline{探究三}{例1 A；变式1 B}
\vskip 0pt minus 1.2pt
% pair:导-课时94-拓展
\ansline{拓展延伸}{光滑表面→电荷密度小→场强弱→不易放电（防漏电损失）；尖端→电荷密集→场强大→易放电（持续放电引雷中和）——同一「尖端密集→强场→易放电」链的两向运用}
\vskip 0pt minus 1.2pt
% pair:导-课时94-评1
% pair:导-课时94-评2
% pair:导-课时94-评3
% pair:导-课时94-评4
% pair:导-课时94-评5
\ansline{评价}{1．A；2．B；3．A；4．大；大（强）；5．静电屏蔽}
\vskip 0pt minus 1.2pt
\columnbreak   % 收尾轮：9.4 尾行溢栏后章末随栏，消单行空页（试验见 _tmpP1成卷0912/_试验-甲页流，6 页五0）
\jietitle{本章易错过关(导学件答案)}
\qufen{过关题组}{5 题全单选（卷面题号 1～5，与件面同号）}
% pair:导-章末-1
\ansitem{1}{C}
\vskip 0pt minus 1.2pt
% pair:导-章末-2
\ansitem{2}{C}
\vskip 0pt minus 1.2pt
% pair:导-章末-3
\ansitem{3}{A}
\vskip 0pt minus 1.2pt
% pair:导-章末-4
\ansitem{4}{D}
\vskip 0pt minus 1.2pt
% pair:导-章末-5
\ansitem{5}{B}
\vskip 0pt minus 1.2pt
\end{multicols}

\clearpage   % 分本切页：甲→乙（multicols 接续页 OTR 溢高根治；四本各起新页，装配轮跨本页码按本切）
% ============================================================
% 【乙】练习件答案（4 节×16 槽；9.2/9.3 第 2 题位＝略位与件面同构）
% ============================================================
\jietitle{9.1\quad 电荷(练习件答案)}
\begin{multicols}{2}
\raggedcolumns
\emergencystretch=1em
\qufen{基础巩固1～10／综合提升11～14／思维探索15～16}{题号与练习件同号}
\setlength{\anshang}{5.8mm}
% pair:练-课时91-1
\ansitem{1}{B}
\vskip 0pt minus 1.2pt
% pair:练-课时91-2
\ansitem{2}{B}
\vskip 0pt minus 1.2pt
% pair:练-课时91-3
\ansitem{3}{A}
\vskip 0pt minus 1.2pt
% pair:练-课时91-4
\ansitem{4}{B}
\vskip 0pt minus 1.2pt
% pair:练-课时91-5
\ansitem{5}{A（选非）}
\vskip 0pt minus 1.2pt
% pair:练-课时91-6
\ansitem{6}{带负电／带正电／张开}
\vskip 0pt minus 1.2pt
% pair:练-课时91-7
\ansitem{7}{A}
\vskip 0pt minus 1.2pt
% pair:练-课时91-8
\ansitem{8}{A}
\vskip 0pt minus 1.2pt
% pair:练-课时91-9
\ansitem{9}{A（$1.76\times10^{11}\,\unit{C/kg}$）}
\vskip 0pt minus 1.2pt
\setlength{\anshang}{7.6mm}   % 题 10 起双位档
% pair:练-课时91-10
\ansitem{10}{A}
\vskip 0pt minus 1.2pt
% pair:练-课时91-11
\ansitem{11}{$-\dfrac{q}{4}$}
\vskip 0pt minus 1.2pt
% pair:练-课时91-12
\ansitem{12}{AC}
\vskip 0pt minus 1.2pt
% pair:练-课时91-13
\ansitem{13}{C}
\vskip 0pt minus 1.2pt
% pair:练-课时91-14
\ansitem{14}{(1) 甲 $+5.0\times10^{-9}\,\unit{C}$、乙 $-5.0\times10^{-9}\,\unit{C}$；(2) P 带负电}
\vskip 0pt minus 1.2pt
% pair:练-课时91-15
\ansitem{15}{(1) 甲带负电、乙带正电，$q_{1}=q_{2}$（等量异号·电荷守恒），自由电子由乙侧向靠近玻璃棒的甲侧移动；(2) $q_{1}=4.0\times10^{-9}\,\unit{C}$；末态甲＝丙＝$-1.0\times10^{-9}\,\unit{C}$、乙＝$+2.0\times10^{-9}\,\unit{C}$；(3) 仍能唯一确定，$q_{1}=4.0\times10^{-9}\,\unit{C}$（$+1.0$ 一支与感应链矛盾舍去）}
\vskip 0pt minus 1.2pt
% pair:练-课时91-16
\ansitem{16}{(1) B＝$-8.0\times10^{-9}\,\unit{C}$；(2) C 可能值＝$+4.0$／$+2.0$／$0$／$-2.0$／$-4.0$（$\times10^{-9}\,\unit{C}$），依序「AC→AB／BC→AC／AB→BC 与 AB→AC／AC→BC／BC→AB」；(3) 不能（三球代数和恒为 $0$）}
\vskip 0pt minus 1.2pt
\end{multicols}
\jietitle{9.2\quad 库仑定律(练习件答案)}
\begin{multicols}{2}
\raggedcolumns
\emergencystretch=1em
\qufen{基础巩固1～10／综合提升11～14／思维探索15～16}{题号与练习件同号}
\setlength{\anshang}{5.8mm}
% pair:练-课时92-1
\ansitem{1}{C（$\dfrac{F}{2}$）}
\vskip 0pt minus 1.2pt
% pair:练-课时92-3
\ansitem{3}{BC}
\vskip 0pt minus 1.2pt
\ansline{说明}{第 2 题略——原中4（源Q16）转导学件探究二例1 照用，练习件让位防双收（批2值台账 §五-1）。}
\vskip 0pt minus 1.2pt
% pair:练-课时92-4
\ansitem{4}{C}
\vskip 0pt minus 1.2pt
% pair:练-课时92-5
\ansitem{5}{B（$F_{2}$）}
\vskip 0pt minus 1.2pt
% pair:练-课时92-6
\ansitem{6}{D}
\vskip 0pt minus 1.2pt
% pair:练-课时92-7
\ansitem{7}{D（$\dfrac{2kQql}{h^{3}}$）}
\vskip 0pt minus 1.2pt
% pair:练-课时92-8
\ansitem{8}{(1) $F_{1}=12\,\unit{N}$；(2) $q_{2}=6\times10^{-5}\,\unit{C}$；(3) $m_{2}=1.6\,\unit{kg}$}
\vskip 0pt minus 1.2pt
% pair:练-课时92-9
\ansitem{9}{C（$\sqrt{\dfrac{m_{1}}{m_{2}}}$）}
\vskip 0pt minus 1.2pt
\setlength{\anshang}{7.6mm}   % 题 10 起双位档
% pair:练-课时92-10
\ansitem{10}{C（$9:8$）}
\vskip 0pt minus 1.2pt
% pair:练-课时92-11
\ansitem{11}{(1) $F=\dfrac{kq^{2}}{2d^{2}}$，方向由 B 指向 A；(2) $L=\dfrac{3(\sqrt{3}-1)d}{2}$}
\vskip 0pt minus 1.2pt
% pair:练-课时92-12
\ansitem{12}{D（$F_{2}=F_{1}\cos^{2}\theta$）}
\vskip 0pt minus 1.2pt
% pair:练-课时92-13
\ansitem{13}{BD}
\vskip 0pt minus 1.2pt
\ansline{解析}{中13 详解侧勘正：错字「出于」→「处于」（用题必改义务，答案值不涉）。}
\vskip 0pt minus 1.2pt
% pair:练-课时92-14
\ansitem{14}{(1) $m=\dfrac{kq^{2}}{gL^{2}}$；(2) $a=2\sqrt{3}g$}
\vskip 0pt minus 1.2pt
% pair:练-课时92-15
\ansitem{15}{C（AC:BC＝$m_{2}:m_{1}$ 不变）}
\vskip 0pt minus 1.2pt
% pair:练-课时92-16
\ansitem{16}{A（$f_{\text{地}}=5\sqrt{3}\,\unit{N}$）}
\vskip 0pt minus 1.2pt
\ansline{解析}{47题11 详解侧摩擦理据已按钉值门抽验重写（$F=5\sqrt{3}\,\unit{N}\to T=10\,\unit{N}\to\mu=\dfrac{\sqrt{3}}{6}$ 四级链全复算在案，批2值台账 §四）；答案值不涉。}
\vskip 0pt minus 1.2pt
\end{multicols}
\jietitle{9.3\quad 电场 电场强度(练习件答案)}
\begin{multicols}{2}
\raggedcolumns
\emergencystretch=1em
\qufen{基础巩固1～10／综合提升11～14／思维探索15～16}{题号与练习件同号}
\setlength{\anshang}{5.8mm}
% pair:练-课时93-1
\ansitem{1}{C（$2:\sqrt{3}$）}
\vskip 0pt minus 1.2pt
% pair:练-课时93-3
\ansitem{3}{D}
\vskip 0pt minus 1.2pt
\ansline{说明}{第 2 题略——原简8（源Q48）本轮出列报缺，不凑题（批3值台账 §五-1）。}
\vskip 0pt minus 1.2pt
% pair:练-课时93-4
\ansitem{4}{C}
\vskip 0pt minus 1.2pt
% pair:练-课时93-5
\ansitem{5}{D（P 负 Q 正）}
\vskip 0pt minus 1.2pt
% pair:练-课时93-6
\ansitem{6}{(1) $F=mg\tan\alpha$；(2) $E_{A}=\dfrac{mg\tan\alpha}{q}$，方向水平向右；(3) $Q=\dfrac{mgr^{2}\tan\alpha}{kq}$}
\vskip 0pt minus 1.2pt
% pair:练-课时93-7
\ansitem{7}{B（$\dfrac{E}{2}$，沿 AO 连线斜向上）}
\vskip 0pt minus 1.2pt
% pair:练-课时93-8
\ansitem{8}{C（$\sqrt{2}E$）}
\vskip 0pt minus 1.2pt
% pair:练-课时93-9
\ansitem{9}{A（$1.0\times10^{3}\,\unit{N/C}$，沿 QP 连线由 Q 指向 P）}
\vskip 0pt minus 1.2pt
\setlength{\anshang}{7.6mm}   % 题 10 起双位档
% pair:练-课时93-10
\ansitem{10}{A（$3.0\times10^{4}\,\unit{N/C}$，与静电力同向）}
\vskip 0pt minus 1.2pt
% pair:练-课时93-11
\ansitem{11}{D（$x_{B}=0.4\,\unit{m}$）}
\vskip 0pt minus 1.2pt
\ansline{解析}{中17 详解侧修复：脱符「xB2」→$x_{B}^{2}$（用题必改义务，答案值不涉）。}
\vskip 0pt minus 1.2pt
% pair:练-课时93-12
\ansitem{12}{B（$\dfrac{7kQ}{36R^{2}}$）}
\vskip 0pt minus 1.2pt
% pair:练-课时93-13
\ansitem{13}{C}
\vskip 0pt minus 1.2pt
% pair:练-课时93-14
\ansitem{14}{B（A 负 B 正，$m_{A}=m_{B}$）}
\vskip 0pt minus 1.2pt
% pair:练-课时93-15
\ansitem{15}{C（$Q_{1}=-2q$，$Q_{2}=+\sqrt{3}q$，$Q_{3}=-2\sqrt{3}q$）}
\vskip 0pt minus 1.2pt
% pair:练-课时93-16
\ansitem{16}{C（$\left|\dfrac{3kq}{8R^{2}}-2E\right|$）}
\vskip 0pt minus 1.2pt
\ansline{解析}{冲4 详解侧已按亲算补全版重写（$\sigma_{+}$ 核算＋镜像 $2E$ 记账，批3值台账 §四 双层割补数值闭合 1.46923687）；答案值不涉。}
\vskip 0pt minus 1.2pt
\end{multicols}
\jietitle{9.4\quad 静电的防止与利用(练习件答案)}
\begin{multicols}{2}
\raggedcolumns
\emergencystretch=1em
\qufen{基础巩固1～10／综合提升11～14／思维探索15～16}{题号与练习件同号}
\setlength{\anshang}{5.8mm}
% pair:练-课时94-1
\ansitem{1}{B}
\vskip 0pt minus 1.2pt
\ansline{解析}{简9 详解侧驳论补全：D 项「避雷线吸收导入大地＝尖端放电」系窄判读——吸收导地≠尖端放电（补全或换 D 为静电复印，两读答案 B 均稳；批4值台账 §六-2②）；答案值不涉。}
\vskip 0pt minus 1.2pt
% pair:练-课时94-2
\ansitem{2}{D（近指端 $+$、远端 $-$）}
\vskip 0pt minus 1.2pt
% pair:练-课时94-3
\ansitem{3}{A}
\vskip 0pt minus 1.2pt
% pair:练-课时94-4
\ansitem{4}{B}
\vskip 0pt minus 1.2pt
% pair:练-课时94-5
\ansitem{5}{垂直；沿导体表面定向移动}
\vskip 0pt minus 1.2pt
% pair:练-课时94-6
\ansitem{6}{处处为零；静电屏蔽；相同}
\vskip 0pt minus 1.2pt
% pair:练-课时94-7
\ansitem{7}{B}
\vskip 0pt minus 1.2pt
% pair:练-课时94-8
\ansitem{8}{B（防尖端＝光滑球形）}
\vskip 0pt minus 1.2pt
% pair:练-课时94-9
\ansitem{9}{A（加湿导走静电）}
\vskip 0pt minus 1.2pt
\setlength{\anshang}{7.6mm}   % 题 10 起双位档
% pair:练-课时94-10
\ansitem{10}{B}
\vskip 0pt minus 1.2pt
% pair:练-课时94-11
\ansitem{11}{B（感应场向左；正解式 $\dfrac{4kq}{(2R+l)^{2}}$ 随详解在册）}
\vskip 0pt minus 1.2pt
\ansline{解析}{中24 详解侧勘正：错字「该次」→「该处」（用题必改义务，答案值不涉）。}
\vskip 0pt minus 1.2pt
% pair:练-课时94-12
\ansitem{12}{B}
\vskip 0pt minus 1.2pt
% pair:练-课时94-13
\ansitem{13}{B}
\vskip 0pt minus 1.2pt
% pair:练-课时94-14
\ansitem{14}{AC}
\vskip 0pt minus 1.2pt
% pair:练-课时94-15
\ansitem{15}{B（A 不偏离＋B 向右偏）}
\vskip 0pt minus 1.2pt
% pair:练-课时94-16
\ansitem{16}{(1) 内表面 $-3.0\times10^{-9}\,\unit{C}$（均匀）、外表面 $+3.0\times10^{-9}\,\unit{C}$（不均匀，背离 $q_{2}$ 侧密）；(2) $E_{A}=300\,\unit{N/C}$，沿 $q_{1}A$ 连线背离 $q_{1}$；$E_{B}=0$；(3) S 闭合或先闭合再断开，外表面均由 $+q_{1}$→约 $0$，壳外不再有由 $q_{1}$ 引起的电场；与一直不接地（外表面 $+q_{1}$、壳外有由 $q_{1}$ 引起的电场）成三支对照}
\vskip 0pt minus 1.2pt
\end{multicols}

\clearpage   % 分本切页：乙→丙
% ============================================================
% 【丙】拓展册答案（拓-001~046 值随批5b 台账 §一 回填；◐改号件按 §一-3 映射；
%   拓46＝收尾轮改单选（主脑回执④·批5b §八-1）；题号与拓展册册面同号）
% ============================================================
\jietitle{拓展册(答案)}
\begin{multicols}{2}
\raggedcolumns
\emergencystretch=1em
\qufen{9.1\quad 电荷}{册面题号 1～2（注记位 1＋单选 1）}
\setlength{\anshang}{5.8mm}
% pair:拓-001
\ansitem{1}{B}
\vskip 0pt minus 1.2pt
\ansline{解析}{注记位（A-2 式）——题面不重收，见 9.1 导学件课中探究·探究点二 例1（中3·源Q10），答案 B 在册（批5b §一-2⑨）。}
\vskip 0pt minus 1.2pt
% pair:拓-002
\ansitem{2}{D}
\vskip 0pt minus 1.2pt
\qufen{9.2\quad 库仑定律}{册面题号 3～19（单选 14＋多选 3）}
% pair:拓-003
\ansitem{3}{C}
\vskip 0pt minus 1.2pt
% pair:拓-004
\ansitem{4}{B}
\vskip 0pt minus 1.2pt
% pair:拓-005
\ansitem{5}{A}
\vskip 0pt minus 1.2pt
% pair:拓-006
\ansitem{6}{C}
\vskip 0pt minus 1.2pt
% pair:拓-007
\ansitem{7}{A}
\vskip 0pt minus 1.2pt
% pair:拓-008
\ansitem{8}{D}
\vskip 0pt minus 1.2pt
% pair:拓-009
\ansitem{9}{A}
\vskip 0pt minus 1.2pt
\setlength{\anshang}{7.6mm}   % 题 10 起双位档
% pair:拓-010
\ansitem{10}{A}
\vskip 0pt minus 1.2pt
% pair:拓-011
\ansitem{11}{D}
\vskip 0pt minus 1.2pt
% pair:拓-012
\ansitem{12}{B}
\vskip 0pt minus 1.2pt
% pair:拓-013
\ansitem{13}{C}
\vskip 0pt minus 1.2pt
% pair:拓-014
\ansitem{14}{D}
\vskip 0pt minus 1.2pt
% pair:拓-015
\ansitem{15}{B}
\vskip 0pt minus 1.2pt
% pair:拓-016
\ansitem{16}{C}
\vskip 0pt minus 1.2pt
% pair:拓-017
\ansitem{17}{CD}
\vskip 0pt minus 1.2pt
% pair:拓-018
\ansitem{18}{BD}
\vskip 0pt minus 1.2pt
% pair:拓-019
\ansitem{19}{AB}
\vskip 0pt minus 1.2pt
\qufen{9.3\quad 电场 电场强度}{册面题号 20～40（单选 17＋多选 3＋填空 1）}
% pair:拓-020
\ansitem{20}{C}
\vskip 0pt minus 1.2pt
% pair:拓-021
\ansitem{21}{D}
\vskip 0pt minus 1.2pt
% pair:拓-022
\ansitem{22}{C}
\vskip 0pt minus 1.2pt
% pair:拓-023
\ansitem{23}{C}
\vskip 0pt minus 1.2pt
% pair:拓-024
\ansitem{24}{D}
\vskip 0pt minus 1.2pt
% pair:拓-025
\ansitem{25}{A}
\vskip 0pt minus 1.2pt
% pair:拓-026
\ansitem{26}{D}
\vskip 0pt minus 1.2pt
% pair:拓-027
\ansitem{27}{B}
\vskip 0pt minus 1.2pt
% pair:拓-028
\ansitem{28}{B}
\vskip 0pt minus 1.2pt
% pair:拓-029
\ansitem{29}{A}
\vskip 0pt minus 1.2pt
% pair:拓-030
\ansitem{30}{A}
\vskip 0pt minus 1.2pt
% pair:拓-031
\ansitem{31}{D}
\vskip 0pt minus 1.2pt
% pair:拓-032
\ansitem{32}{A（合场强为 $0$）}
\vskip 0pt minus 1.2pt
% pair:拓-033
\ansitem{33}{B}
\vskip 0pt minus 1.2pt
% pair:拓-034
\ansitem{34}{B（$\dfrac{3kQ}{l_{0}^{2}}$）}
\vskip 0pt minus 1.2pt
% pair:拓-035
\ansitem{35}{B}
\vskip 0pt minus 1.2pt
% pair:拓-036
\ansitem{36}{B（$\dfrac{45kq}{16h^{2}}$）}
\vskip 0pt minus 1.2pt
% pair:拓-037
\ansitem{37}{BC}
\vskip 0pt minus 1.2pt
% pair:拓-038
\ansitem{38}{BC}
\vskip 0pt minus 1.2pt
% pair:拓-039
\ansitem{39}{BC}
\vskip 0pt minus 1.2pt
% pair:拓-040
\ansitem{40}{$<$；$\unit{kg\cdot m^{3}\cdot s^{-4}\cdot A^{-2}}$}
\vskip 0pt minus 1.2pt
\qufen{9.4\quad 静电的防止与利用}{册面题号 41～46（单选 4＋多选 2 题位；拓46 收尾轮改单选、组行不动）}
% pair:拓-041
\ansitem{41}{D}
\vskip 0pt minus 1.2pt
% pair:拓-042
\ansitem{42}{A}
\vskip 0pt minus 1.2pt
% pair:拓-043
\ansitem{43}{D}
\vskip 0pt minus 1.2pt
% pair:拓-044
\ansitem{44}{B}
\vskip 0pt minus 1.2pt
% pair:拓-045
\ansitem{45}{CD}
\vskip 0pt minus 1.2pt
% pair:拓-046
\ansitem{46}{D（$b$、$c$ 两点的电场强度均为零）}
\vskip 0pt minus 1.2pt
\ansline{解析}{收尾轮改单选（主脑回执④）：正解取 49题9 详解正确结论（原B取反式）；干扰项错因——A 系「腔内无场线$\Rightarrow$感应场亦为零」半截子理解（合场零恰证 $E_{\text{感}}(c)=-E_{0}\neq0$）、B 违静电平衡表面场垂直性（壳面切向分量恒零）、C 违极向聚集增强（$a$ 在场轴：$E_{a}=E_{0}(1+2R^{3}/x^{3})>E_{0}$）；题号/源号/组行不动（批5b §八-1）。}
\vskip 0pt minus 1.2pt
\end{multicols}

\clearpage   % 分本切页：丙→丁
% ============================================================
% 【丁】测评卷答案（测-1~19 值随批5b 台账 §四 回填；卷面题号 1～19 同号；卷末速查对账＝收尾轮 §八-5/#16）
% ============================================================
\jietitle{单元素养测评卷(一)(答案)}
\begin{multicols}{2}
\raggedcolumns
\emergencystretch=1em
\qufen{一、选择题}{本题共 8 小题，每小题 4 分，共 32 分（卷面题号 1～8）}
\setlength{\anshang}{5.8mm}
% pair:测-1
\ansitem{1}{B}
\vskip 0pt minus 1.2pt
% pair:测-2
\ansitem{2}{A}
\vskip 0pt minus 1.2pt
% pair:测-3
\ansitem{3}{C}
\vskip 0pt minus 1.2pt
% pair:测-4
\ansitem{4}{D}
\vskip 0pt minus 1.2pt
% pair:测-5
\ansitem{5}{C}
\vskip 0pt minus 1.2pt
% pair:测-6
\ansitem{6}{D}
\vskip 0pt minus 1.2pt
% pair:测-7
\ansitem{7}{B}
\vskip 0pt minus 1.2pt
% pair:测-8
\ansitem{8}{D}
\vskip 0pt minus 1.2pt
\qufen{二、选择题}{本题共 3 小题，每小题 6 分，共 18 分（卷面题号 9～11）}
% pair:测-9
\ansitem{9}{AC}
\vskip 0pt minus 1.2pt
\setlength{\anshang}{7.6mm}   % 题 10 起双位档
% pair:测-10
\ansitem{10}{CD}
\vskip 0pt minus 1.2pt
% pair:测-11
\ansitem{11}{AC}
\vskip 0pt minus 1.2pt
\qufen{三、填空题}{本题共 3 小题，每小题 5 分，共 15 分（卷面题号 12～14）}
% pair:测-12
\ansitem{12}{带负电；带正电；张开}
\vskip 0pt minus 1.2pt
% pair:测-13
\ansitem{13}{$<$；$\unit{kg\cdot m^{3}\cdot s^{-4}\cdot A^{-2}}$}
\vskip 0pt minus 1.2pt
% pair:测-14
\ansitem{14}{处处为零；$\dfrac{4kq}{(2R+l)^{2}}$；水平向左（指向 $q$ 一侧）}
\vskip 0pt minus 1.2pt
\qufen{四、解答题}{本题共 5 小题，共 35 分（卷面题号 15～19）}
% pair:测-15
\ansitem{15}{$-\dfrac{q}{4}$}
\vskip 0pt minus 1.2pt
% pair:测-16
\ansitem{16}{$mg\tan\alpha$；$\dfrac{mg\tan\alpha}{q}$，水平向右；$\dfrac{mgr^{2}\tan\alpha}{kq}$}
\vskip 0pt minus 1.2pt
% pair:测-17
\ansitem{17}{$12\,\unit{N}$；$6\times10^{-5}\,\unit{C}$；$1.6\,\unit{kg}$}
\vskip 0pt minus 1.2pt
% pair:测-18
\ansitem{18}{$\dfrac{kq^{2}}{gL^{2}}$；$2\sqrt{3}g$}
\vskip 0pt minus 1.2pt
% pair:测-19
\ansitem{19}{$\left|\dfrac{3kq}{8R^{2}}-2E\right|$}
\vskip 0pt minus 1.2pt
\end{multicols}
